\chapter{Bağlantılılık}

Bağlantılılık, bir topolojik uzayın iki ayrı açık parçaya bölünememesi özelliğidir.
Yol-bağlantılılık daha güçlü bir kavramdır ve sürekli yolların varlığına dayanır.

%-------------------------------------------------------------------
\section{Konu}

\subsection{Bağlantılılık Kavramları}

\begin{definition}[Bağlantılı Uzay]
$(X, \tau)$ uzayı \textbf{bağlantılı} ise: $X = U \cup V$, $U \cap V = \emptyset$,
$U,V$ açık $\Rightarrow U = \emptyset$ veya $V = \emptyset$.
\end{definition}

\begin{definition}[Yol-Bağlantılı]
$\forall x,y \in X$: $\exists$ sürekli $f:[0,1] \to X$, $f(0) = x$, $f(1) = y$.
\end{definition}

\begin{table}[h]
\centering
\begin{tabular}{ll}
\toprule
\textbf{Kavram} & \textbf{Özellik} \\
\midrule
Bağlantılı & Clopen bölme yok \\
Yol-bağlantılı & Her iki nokta arasında sürekli yol var \\
Ark-bağlantılı & Yol-bağlantılı; yollar enjektif \\
Lokal bağlantılı & Her noktanın bağlantılı komşuluğu \\
Tamamen bağlantısız & Her bağlantılı alt küme tek noktadır \\
\bottomrule
\end{tabular}
\caption{Bağlantılılık hiyerarşisi}
\end{table}

\textbf{Sıralama:} Ark-bağlantılı $\Rightarrow$ Yol-bağlantılı $\Rightarrow$ Bağlantılı.

\subsection{Clopen Ayrılma}

$X$ bağlantısız $\Leftrightarrow$ boş olmayan trivial olmayan bir clopen $A \subsetneq X$ vardır.

%-------------------------------------------------------------------
\section{Teoremler}

\begin{theorem}
Bağlantılı küme süreklilik altında bağlantılıdır:
$f: X \to Y$ sürekli, $X$ bağlantılı $\Rightarrow f(X)$ bağlantılı.
\end{theorem}

\begin{theorem}
Yol-bağlantılı $\Rightarrow$ bağlantılı. Tersi genel olarak yanlış.
(Karşı örnek: Topolojist sinüs eğrisi.)
\end{theorem}

\begin{theorem}[Ara Değer Teoremi]
$f: X \to \mathbb{R}$ sürekli, $X$ bağlantılı $\Rightarrow f(X)$ bir aralıktır.
\end{theorem}

\begin{theorem}
Gerçek doğru $\mathbb{R}$'de bağlantılı alt kümeler tam olarak aralıklardır.
\end{theorem}

%-------------------------------------------------------------------
\section{Algoritmalar}

\subsection{Sonlu Bağlantılılık — $O(|\tau|^2)$}

\begin{algorithm}
\caption{BagliMi$(X, \tau)$}
\begin{algorithmic}[1]
\For{each non-empty $A \subsetneq X$}
    \If{$A \in \tau$ \textbf{and} $X \setminus A \in \tau$}
        \State \Return False \Comment{Clopen bölme bulundu}
    \EndIf
\EndFor
\State \Return True
\end{algorithmic}
\end{algorithm}

Karmaşıklık: $O(|\tau|^2)$ — clopen küme taraması.

%-------------------------------------------------------------------
\section{pytop API}

\begin{lstlisting}[language=Python]
from pytop import (
    is_connected,
    is_path_connected,
    is_locally_connected,
    is_totally_disconnected,
)
\end{lstlisting}

%-------------------------------------------------------------------
\section{Örnekler}

\subsection*{Örnek 5.1 — Sierpiński: Bağlantılı}

\begin{lstlisting}[language=Python]
s = sierpinski_space()
print("connected?", is_connected(s).status)
print("path_connected?", is_path_connected(s).status)
\end{lstlisting}

\begin{verbatim}
connected?       true
path_connected?  unknown
\end{verbatim}

\subsection*{Örnek 5.2 — Ayrık Topoloji: Tamamen Bağlantısız}

\begin{lstlisting}[language=Python]
d = discrete_topology(1, 2, 3)
print("connected?", is_connected(d).status)
print("totally_disconn?", is_totally_disconnected(d).status)
\end{lstlisting}

\begin{verbatim}
connected?          false
totally_disconn?    true
\end{verbatim}

\subsection*{Örnek 5.3 — Gerçek Doğru ve [0,1]}

\begin{lstlisting}[language=Python]
rl = real_line_metric()
ui = closed_unit_interval_metric()
print("R connected?", is_connected(rl).status)
print("[0,1] path_connected?", is_path_connected(ui).status)
\end{lstlisting}

\begin{verbatim}
R connected?          true
[0,1] path_connected? true
\end{verbatim}

%-------------------------------------------------------------------
\section{Alıştırmalar}

\subsection*{Kodlama}

\begin{enumerate}[label=K\arabic*.]
\item \texttt{make\_topology(\{1,2,3\},\{1\},\{2,3\})} için \texttt{is\_connected} hesaplayın.
\item \texttt{finite\_chain\_space(4)} zincirinin bağlantılı olup olmadığını kontrol edin.
\item İki noktalı ayrık ve indiscrete uzaylar üzerinde \texttt{is\_connected},
      \texttt{is\_path\_connected} karşılaştırın.
\end{enumerate}

\subsection*{Teori}

\begin{enumerate}[label=T\arabic*.]
\item Bağlantılı + sürekli $\Rightarrow$ görüntü bağlantılı teoremini ispatlayın.
\item Yol-bağlantılı $\Rightarrow$ bağlantılı implicasyonunu ispatlayın.
\end{enumerate}
